% Test des interactions de bout en bout d'une promotion de 1 000 étudiants, par l'interface web.
% Compilation : latexmk -pdf test-des-interactions.tex
\documentclass[11pt,a4paper]{article}
\usepackage[french]{babel}
\usepackage[T1]{fontenc}
\usepackage[utf8]{inputenc}
\usepackage{lmodern,microtype,booktabs,tabularx,xcolor,hyperref,array,amsmath}
\usepackage[margin=2.2cm]{geometry}
\hypersetup{colorlinks=true, linkcolor=blue!50!black, urlcolor=blue!50!black}
\newcommand{\code}[1]{\texttt{\detokenize{#1}}}
\newcolumntype{R}{>{\raggedleft\arraybackslash}X}
\title{Test des interactions à 1\,000 étudiants}
\author{Sébastien Lefebvre}
\date{Novembre 2026 -- versions 1.0.3 et 1.0.5 du module}

\begin{document}
\maketitle

\section{Objet et méthode}

Les tests précédents mesuraient des pages (test de charge), des courses à la même ressource
(test de concurrence) et la tenue d'un serveur saturé (test de stress). Celui-ci fait vivre une
promotion entière, \textbf{de la première inscription au supplément au diplôme émis}, par
l'interface web seule, comme des navigateurs : chaque formulaire est lu dans la page, avec ses
champs cachés et sa clé de session ; chaque justificatif est déposé dans la zone brouillon par
le sélecteur de fichiers de Moodle. Puis il vérifie en base que chaque action a laissé exactement
la trace attendue.

\paragraph{Serveur.} Le banc habituel : une machine de 4 c\oe urs et 15~Go, nginx, PHP-FPM~8.3
(16 processus, OPcache, JIT), PostgreSQL~16, Redis pour les sessions, APCu pour le cache
applicatif, Moodle~4.5. Le client tourne sur la même machine.

\paragraph{Promotion.} Une promotion \textbf{vierge} (\code{VIERGE=1 charge_setup.php}) :
1\,000 étudiants inscrits au cours, chacun avec son enseignant référent dans « Gestion des
stages » ; 150 enseignants ; 4 membres de la DEVE ; 60 EP publiés au catalogue (un sur quatre noté
sur 20, seuil 10 ; places limitées à 40 pour deux EP sur trois), avec un ou deux responsables
chacun ; minimum de 6 ECTS en troisième année ; circuit des éléments du supplément
« référent puis DEVE ». Aucun crédit, aucun élément, aucune décision au départ : tout ce qui est
compté à la fin est né des formulaires. Les enseignants et la DEVE sont ceux des six promotions du
test de stress, et la synthèse enseignant est propre à la promotion.

\paragraph{Scénario.} \code{docs/audit/banc/interactions.py}, 40 acteurs en parallèle dans chaque
phase :
\begin{enumerate}
  \item \textbf{Étudiants} (1\,000) : tableau de bord ; une page du catalogue ; trois inscriptions
    directes depuis la liste ; une déclaration hors catalogue avec justificatif PDF ; une demande
    d'élément du supplément avec preuve PDF et période ; un sur dix retire ensuite sa déclaration,
    un sur dix la modifie.
  \item \textbf{Responsables d'EP} (55, ceux des EP qui ont des inscrits) : notes ou validations
    de toute la liste des participants en un envoi (une note sur quatre sous le seuil, une
    décision sur dix négative) ; un EP noté par import CSV (prévisualisation puis confirmation).
  \item \textbf{Référents} (150) : chaque déclaration de leurs étudiants, ouverte puis validée
    (neuf sur dix) ou refusée avec un commentaire ; puis chaque élément du supplément, tranché
    depuis la page « À traiter » de la synthèse.
  \item \textbf{DEVE} : paramétrage du supplément (textes des rubriques, signataire) ; décision
    finale sur les 1\,000 éléments, les quatre comptes de la DEVE se partageant la file ;
    validation des années en lot (atteintes, puis non atteintes) ; émission groupée des 1\,000
    suppléments (tâche de fond) ; vérification publique de 50 suppléments, et de 50 faux codes ;
    exports CSV et export complet.
  \item \textbf{Contrôles en base}.
\end{enumerate}

\section{Déroulement}

\begin{center}
\begin{tabularx}{\linewidth}{@{}X r r r r@{}}
\toprule
\textbf{Phase} & \textbf{Acteurs} & \textbf{Durée} & \textbf{Requêtes} & \textbf{Requêtes/s} \\
\midrule
1. Étudiants : inscriptions, déclarations, éléments & 1\,000 & 254,7 s & 13\,500 & 53,0 \\
2. Responsables : notes et validations en lot & 55 & 23,7 s & 222 & 9,4 \\
3. Référents : déclarations, éléments (synthèse) & 150 & 77,2 s & 3\,685 & 47,7 \\
4. DEVE : paramétrage du supplément & 1 & 0,4 s & 4 & -- \\
4a. DEVE : décision finale sur les éléments & 4 & 42,8 s & 1\,115 & 26,0 \\
4b. DEVE : validation des années en lot & 1 & 1,4 s & 5 & -- \\
4c. DEVE : émission groupée des suppléments & 1 & 155,1 s & 4 & -- \\
4d. Public : vérification de 50 suppléments et 50 faux codes & 50 & 0,6 s & 100 & 173,9 \\
4e. DEVE : exports CSV et export complet & 1 & 6,4 s & 5 & -- \\
\midrule
\textbf{Total} & & \textbf{9 min 22 s} & \textbf{18\,640} & \\
\bottomrule
\end{tabularx}
\end{center}

\begin{itemize}
  \item Un étudiant fait tout son parcours (14 pages et envois, deux dépôts de fichiers) en
    10~s en médiane, avec 39 autres en même temps : la promotion entière passe en un peu plus de
    4 minutes, au rythme de 53 requêtes par seconde -- pages à formulaire, dépôts et envois
    compris, soit un peu moins que les 60 pages par seconde de simple consultation du test de
    stress.
  \item L'enregistrement d'une liste de participants (40 à 50 notes, avec un message à chaque
    étudiant dont la décision est finale) prend 3 à 4~s.
  \item L'émission groupée produit les 1\,000 suppléments, chacun avec son PDF, son empreinte et
    son sceau, en 140~s de tâche de fond (0,14~s et 25 requêtes à la base par supplément).
  \item L'export complet de la promotion pèse 133~Mo (classeur, paramétrage, données, et les
    3\,000 fichiers : justificatifs, preuves, suppléments) et se télécharge en 6~s ; les exports
    CSV en moins d'une seconde.
\end{itemize}

\section{Contrôles}

\begin{center}
\begin{tabularx}{\linewidth}{@{}X r r@{}}
\toprule
\textbf{Contrôle} & \textbf{Attendu} & \textbf{Constaté} \\
\midrule
Inscriptions faites depuis le catalogue & 2\,984 & 2\,984 \\
\quad dont validées / refusées par les responsables & & 2\,570 / 414 \\
Déclarations, chacune avec son justificatif & 1\,000 & 1\,000 \\
\quad dont retirées par l'étudiant & 100 & 100 \\
\quad dont modifiées par l'étudiant avant décision & 100 & 100 \\
\quad dont validées / refusées par le référent & & 811 / 89 \\
Demandes encore en attente & 0 & 0 \\
Demandes tranchées sans décision dans l'historique & 0 & 0 \\
Décisions enregistrées deux fois pour la même étape & 0 & 0 \\
Inscriptions en double, places dépassées & 0 & 0 \\
Éléments du supplément, chacun avec sa preuve & 1\,000 & 1\,000 \\
\quad validés par le référent puis par la DEVE & 1\,000 & 1\,000 \\
Décisions d'année (3\textsuperscript{e} année) & 1\,000 & 1\,000 \\
\quad années validées (minimum de 6 ECTS atteint) / non validées & & 605 / 395 \\
Suppléments émis, scellés, avec leur PDF & 1\,000 & 1\,000 \\
Numéros de supplément en double & 0 & 0 \\
Vérifications publiques acceptées / faux codes refusés & 50 / 50 & 50 / 50 \\
Verrous restés pris en fin de test & 0 & 0 \\
Anomalies relevées dans les pages (erreur, avertissement PHP, message d'erreur) & 0 & 0 \\
\bottomrule
\end{tabularx}
\end{center}

Les refus sont ceux du scénario : notes de 8 sur 20 (sous le seuil) pour un quart des inscrits
aux EP notés, une décision négative sur dix ailleurs. Les 395 années non validées sont celles des
étudiants dont les ECTS validés n'atteignent pas 6, à cause de ces refus : la validation en lot a
validé les autres et marqué celles-là, sans en laisser aucune sans décision.

\section{Défauts trouvés et corrigés}

\subsection{Connexions refusées après un redémarrage de PHP-FPM}

Au premier passage, \textbf{8 étudiants sur 1\,000} n'ont rien pu faire : leur connexion
aboutissait à une erreur 500 de la page d'accueil de Moodle (\code{/my/}), toutes dans la même
seconde, au tout début. Le second passage n'en a eu aucune. La différence : PHP-FPM venait d'être
redémarré avant le premier.

L'erreur, reproduite à volonté, est \og Unable to acquire a lock for caching\fg{}, levée par la
fonction \code{build_course_cache()} de \code{course_modinfo}, c'est-à-dire par le c\oe ur de Moodle quand il
reconstruit le cache des activités d'un cours. Le cache applicatif est dans APCu (recommandation
du test de stress), qui vit dans la mémoire de PHP-FPM : un redémarrage, comme une purge des
caches, le vide. APCu n'a pas de verrouillage natif ; Moodle~4.5 prend alors un verrou de fichier
\emph{sans attendre} (la méthode \code{lock()} de \code{cachelock_file} est appelée sans
blocage) : toute requête qui veut
reconstruire le cache d'un cours pendant qu'une autre le fait échoue aussitôt, en 0,8~s. Le
module n'y est pour rien (il ne fournit pas d'informations au cache des cours, dont la
reconstruction prend 5 à 21~ms), mais la recommandation d'APCu l'expose.

\begin{center}
\begin{tabularx}{\linewidth}{@{}X r@{}}
\toprule
\textbf{40 premières connexions simultanées} & \textbf{Échecs} \\
\midrule
Tout le cache applicatif dans APCu, après purge des caches (deux essais) & 12, 14 \\
Idem, après un simple redémarrage de PHP-FPM (vague de 60 étudiants) & 14 sur 60 \\
Idem, attente du verrou de fichier portée de 100 à 600 essais & 4 sur 60 \\
\code{core/coursemodinfo} sur le magasin fichier, APCu pour le reste (trois essais) & 0, 0, 0 \\
\bottomrule
\end{tabularx}
\end{center}

Allonger l'attente ne change rien, puisque Moodle n'attend pas. Le remède est d'associer la seule
définition \code{core/coursemodinfo} au magasin fichier par défaut, qui verrouille nativement et
attend (\code{docs/audit/banc/muc_modinfo.php}). Le débit n'en souffre pas : sur un palier de
100 utilisateurs du test de stress, 56,0 et 55,4 pages par seconde avec ce réglage, 50,7 sans
(le premier essai sans ce réglage s'est d'ailleurs arrêté sur ce même échec de connexion).
\code{INSTALL.md} le prescrit désormais avec APCu.

\subsection{Entité HTML affichée en clair}

Les captures du guide ont montré, sur la page des participants d'un EP noté, \og Seuil de
validation\&nbsp;: 10 / 20\fg{} : le séparateur de libellé de la langue française
(\code{&nbsp;:}) passait par \code{s()} et était échappé une seconde fois. Même défaut dans les
dates de la feuille d'émargement et dans la colonne de résultat de la prévisualisation d'un
import de notes. Ces trois libellés emploient désormais le séparateur en texte brut.

\subsection{Titre en double}

La page publique de vérification d'un supplément affichait son titre deux fois (titre de page de
Moodle et titre du contenu). Le second est retiré.

\section{Liaison avec « Gestion des stages »}

Le scénario précédent laissait le barème du type stage à 0 : aucun EP de stage n'y était
attribué. La liaison avec \code{mod_stage} a donc été vérifiée à part, par les pages des deux
modules (\code{liaison_stage.py}), sur la même promotion : 300 stages préparés par les fonctions
de \code{mod_stage} (200 complémentaires, dont un sur sept à l'étranger, 50 obligatoires,
50 complémentaires laissés à valider).

\paragraph{Ce que lit le module.} Quatre tables de \code{mod_stage} (saisies, détail de convention,
enseignants référents, thématiques), en lecture seule ; aucun champ ne manque dans la version
courante de \code{mod_stage}, et l'état « validé DEVE » y est bien l'état final d'un stage.

\paragraph{Défaut trouvé.} Les EP de stage n'étaient recalculés qu'à l'ouverture de trois pages
du module, au plus toutes les cinq minutes, et la nuit. Deux conséquences : un stage que la DEVE
venait de valider dans « Gestion des stages » ne comptait pas si elle passait aussitôt à la
\textbf{validation des années} (année refusée à tort) ou à l'\textbf{émission d'un supplément}
(supplément certifié sans ce stage) ; et une durée retenue corrigée après validation restait
fausse jusqu'à cinq minutes sur les autres pages. \code{mod_stage} ne déclenche pas d'événement
pour une telle correction : écouter ses événements n'aurait pas suffi. Six tests unitaires
reproduisaient le défaut.

\paragraph{Correctif (1.0.5).} Une empreinte des stages complémentaires validés (nombre,
identifiants, durées retenues, années, mobilité, dates de modification) et du barème est comparée
à celle du dernier calcul ; le recalcul a lieu dès qu'elle change, à l'ouverture des pages qui
décomptent, et la validation des années et l'émission d'un supplément l'attendent.

\begin{center}
\begin{tabularx}{\linewidth}{@{}X l@{}}
\toprule
\textbf{Contrôle, après le correctif} & \textbf{Résultat} \\
\midrule
Barème de 0,5 ECTS par jour saisi dans le formulaire du type & conforme \\
250 stages validés en masse dans la page DEVE de \code{mod_stage} & 0,5 s \\
Premier tableau de bord ouvert : 200 crédits, $\text{durée} \times 0{,}5$ chacun & 0,29 s \\
Stages obligatoires et non validés : aucun crédit & conforme \\
20 durées revalidées (page DEVE) et 20 corrigées dans la saisie (sans événement) & suivies \\
10 stages annulés, 10 repassés en obligatoire & crédits retirés \\
Stage validé puis, aussitôt, décision d'année & le stage compte \\
Stage validé puis, aussitôt, émission du supplément & le stage figure au PDF \\
Référent changé dans \code{mod_stage} & la déclaration en attente passe au nouveau \\
Tâche de la nuit ensuite & rien à recalculer \\
Coût de l'empreinte / de la vérification à l'ouverture d'une page & 1,4 ms / 1,7 ms \\
\bottomrule
\end{tabularx}
\end{center}

\section{Ce que le test ne couvre pas}

\begin{itemize}
  \item Les propositions d'EP et les syllabus, les imports du catalogue et des crédits, la
    rentrée et le transfert d'un étudiant, la restauration d'un export : ils ont leurs
    procédures de bout en bout (rapport \emph{Test de charge}) et leurs tests unitaires.
  \item L'envoi réel des courriels : les messages sont produits, le banc ne les expédie pas.
  \item Les navigateurs eux-mêmes (JavaScript, rendu) : les captures du guide utilisateur, prises
    avec Chromium sur cette promotion après le test, en tiennent lieu pour l'affichage.
\end{itemize}

\section{Reproduire}

{\small
\begin{verbatim}
cd docs/audit/banc
VIERGE=1 PERSONNEL=a sudo -E -u www-data php charge_setup.php x 1000 150 60 0
sudo -u www-data php synthese_setup.php site-x.json      # affiche le cmid de la synthèse
python3 interactions.py site-x.json <cmid> > interactions.json
\end{verbatim}
}

Le résumé est écrit sur la sortie d'erreur, le détail (durées, anomalies, contrôles) en JSON sur
la sortie standard. Les variables \code{PHASES}, \code{DEBUT}, \code{LIMITE} et \code{REFERENTS}
rejouent une partie du scénario : c'est ainsi qu'une seconde vague d'étudiants, laissée en
attente, a servi aux captures du guide utilisateur.

\end{document}
